\section{Continuous Governed Business Improvement}
\label{sec:aion-governed-business-improvement}

AION is intended to learn how a particular business operates, but learning cannot be allowed to
become an invisible expansion of authority.  The continuous-improvement layer therefore separates
what a person said, what they corrected, what was actually attempted, what an independent authority
verified, and what change may safely be promoted.  A model response and a conversation transcript
are not outcome evidence.

\subsection{Separate Evidence Rails}

Owner corrections and verified outcomes use different owner-controlled ledgers.  A correction
records its identifier, task class, target reference, actor, corrected-value hash and optional
conversation-reference hash.  It does not copy the conversation or grant authority.  A verified
outcome records the route, procedure, score, observation time and independent authority reference.
This prevents a frequently repeated statement from being mistaken for repeated real-world success.

\subsection{Task-Specific, Time-Bounded Preferences}

Route and procedure performance is grouped by task class.  Invoice matching, sales forecasting and
customer-support triage therefore learn independently.  A route becomes preferred only after at
least three fresh verified outcomes.  Candidates are ranked by measured score, evidence count and
recency.  Every preference has a validity limit; after that time it returns
\emph{stale revalidation required} and supplies no preferred route until refreshed by evidence.

This distinction is important when a provider changes a model, an application changes its interface,
or the business changes its process.  Yesterday's successful route is useful evidence, but it is not
permanent authority.

\subsection{Automation Is Proposed, Not Assumed}

AION may propose automation only after repeated high-quality verified outcomes from the same
preferred route and procedure.  Mixed success across unrelated routes does not satisfy the gate.
The proposal binds the task class, route, procedure and evidence hashes into one proposal hash.
The owner must approve that exact proposal.  Even after approval its state is
\emph{approved but unactivated}: execution authority must still come from the relevant capability
and organisational policy.  This keeps behavioural convenience separate from permission.

\subsection{Licensed and Consented Dataset Curation}

The dataset registry accepts only approved licences such as customer-owned, consented internal,
CC0, CC BY 4.0, Apache 2.0 or MIT material.  A consent reference and at least one permitted purpose
are required.  Fields resembling passwords, keys, tokens, payment-card data, salary or medical data
are rejected until redacted upstream.  The registry retains source and consent hashes, record count,
content hash, licence and purposes; raw records are not copied into the registry.  Revocation makes
the dataset ineligible for new releases.

\subsection{Reviewed and Reversible Model Releases}

Fine-tuning, distillation and learned adapters are release processes, not background mutations.
Each candidate binds a release identifier, base model, eligible dataset identifiers, training method,
human reviewer and rollback target.  It begins inactive and requires evaluation.

Promotion requires both a passing retention cohort after a real minimum elapsed period and a passing
unfamiliar-transfer cohort whose evidence is source-disjoint from training.  A failed required cohort
blocks promotion.  When promoted, the prior rollback target is retained and the change is entered in
the customer change view.  This mechanism can govern local, customer-cloud or replaceable provider
models without making any model the owner of AION's memory or authority.

\subsection{Customer Inspection and Reversal}

The customer change view shows the change type, accountable actor, time, reason, detail hash,
reversibility and exact reverse action.  Supported reversals can disable an approved automation or
roll a model release back to its declared predecessor.  Reversal is itself attributed to the owner
and cannot be replayed silently.

\subsection{Qualification Boundary}

Automated tests verify correction/outcome separation, fresh-evidence thresholds, stale preference
expiry, exact automation approval, dataset licensing and redaction, revocation, reviewed releases,
elapsed retention, source-disjoint transfer and customer reversal.  This layer governs whether a
change may be admitted.  It does not claim to provide a production GPU training service, create
licensed datasets automatically or replace independent customer-domain evaluation.  Those artefacts
and authorities must be connected and qualified in the customer's own environment.

The result is a business brain that can improve continuously while remaining inspectable, reversible
and owned by the customer.
